\section{\texorpdfstring{A uniform lower location for
$\chi$}{A uniform lower location for chi}}
\label{a-valid-unrestricted-lower-location-for-chi}

Let $r_+(n)=r_{S_+,0}$ be the unrestricted zero from Lemma~3.1, and put
\begin{equation*}
  k_\chi^-:=\lfloor r_+(n)\rfloor-\lceil\log n\rceil.
  \tag{4.1}
\end{equation*}
This is a deterministic integer. Since
$r_+(n)=\Theta(n/\log n)$ uniformly in the phase, one has
$1\le k_\chi^-<n$ for all sufficiently large $n$.

For an integer profile
$\mathbf k=(k_i)_{-1\le i\le\alpha-1}$, let
$X_{\mathbf k}$ be the number of unordered proper colorings having exactly
$k_i$ classes of size $\alpha-i$. Direct enumeration gives
\begin{equation*}
  \Exp{X_{\mathbf k}}
  =\frac{n!}{\prod_i((\alpha-i)!)^{k_i}k_i!}
   2^{-\sum_i k_i\binom{\alpha-i}{2}}.
  \tag{4.2}
\end{equation*}
For an integer $k$, let $\mathcal K_{n,k,\alpha+1}$ be the set of these
profiles satisfying
\[
  \sum_i k_i=k,
  \qquad
  \sum_i(\alpha-i)k_i=n,
\]
and define
\[
  E_{n,k,\alpha+1}
  :=\sum_{\mathbf k\in\mathcal K_{n,k,\alpha+1}}
    \Exp{X_{\mathbf k}}.
\]
Thus $E_{n,k,\alpha+1}$ is the expected number of unordered proper colorings
with exactly $k$ nonempty parts, each of size at most $\alpha+1$.

\subsection*{The profile sum}

There are $\alpha+1=O(\log n)$ profile coordinates, and each coordinate lies
between zero and $n$. Hence
\begin{equation*}
  |\mathcal K_{n,k,\alpha+1}|
  \le(n+1)^{\alpha+1}
  =\exp\{O((\log n)^2)\}
  \tag{4.3a}
\end{equation*}
uniformly in $k$ and in the phase. Applying the uniform factorial estimate
(1.2) to (4.2), with the convention $0\log0=0$, gives
\[
  \log\Exp{X_{\mathbf k}}
  \le L_{\mathbf k}+O((\log n)^2).
\]
The error is uniform even when some $k_i=0$: only $O(\log n)$ factorials
occur, and every one contributes $O(\log n)$. Maximizing over the feasible
profile and then summing (4.3a) yields
\begin{equation*}
  \log E_{n,k,\alpha+1}
  \le L_+(n,k)+O((\log n)^2).
  \tag{4.3}
\end{equation*}

\subsection*{Displacement below the root}

Set
\[
  \Delta_n:=r_+(n)-k_\chi^-.
\]
The floor and ceiling in (4.1) give the deterministic bounds
\begin{equation*}
  \log n\le\Delta_n<\log n+2.
  \tag{4.3b}
\end{equation*}
We next verify that the entire interval
$[k_\chi^-,r_+(n)]$ lies in the derivative corridor of Lemma~3.1. Write
$s(k)=n/k$. Uniformly on this interval,
$k=\Theta(n/\log n)$ and $s(k)=\Theta(\log n)$. Therefore
\[
  |s(k_\chi^-)-s(r_+)|
  \le
  \sup_{k\in[k_\chi^-,r_+]}
  \frac{n}{k^2}\,\Delta_n
  =O\!\left(\frac{(\log n)^3}{n}\right)
  =o\!\left(\frac{\log\log n}{\log n}\right).
\]
Thus, after one phase-independent threshold, every point of the interval lies
inside the common corridor used in (3.7).

Equation (3.7), with $S=S_+$ and $c=0$, consequently gives an absolute
$c_*>0$ such that
\begin{equation*}
  \frac{d}{dk}L_+(n,k)\ge c_*(\log n)^2
  \qquad
  (k_\chi^-\le k\le r_+)
  \tag{4.3c}
\end{equation*}
for all sufficiently large $n$, uniformly in the phase. Since
$L_+(n,r_+)=0$, the mean-value theorem and (4.3b) imply
\begin{equation*}
  L_+(n,k_\chi^-)
  \le-c_*(\log n)^2\Delta_n
  \le-c_*(\log n)^3.
  \tag{4.4}
\end{equation*}
Combining (4.3) and (4.4), and decreasing the constant once, gives a
deterministic sequence
\[
  \varepsilon_n^{\mathrm{prof}}
  :=
  \exp\{-c_\chi(\log n)^3\}
  \longrightarrow0
\]
such that
\begin{equation*}
  E_{n,k_\chi^-,\alpha+1}
  \le\varepsilon_n^{\mathrm{prof}}.
  \tag{4.4a}
\end{equation*}

\subsection*{Removing the size cap}

Let
\[
  \mathcal A_n:=\{\alpha(G_n)\le\alpha+1\}.
\]
Equation (2.9) gives a deterministic phase-uniform sequence
$\varepsilon_n^{\mathrm{cap}}\to0$ such that
\[
  \Prob{\mathcal A_n^c}
  \le\varepsilon_n^{\mathrm{cap}}.
\]
On $\mathcal A_n$, every class of every proper coloring has size at most
$\alpha+1$.

Suppose that $\chi(G_n)\le k_\chi^-$ and choose a proper coloring with
$h\le k_\chi^-$ nonempty classes. If $h<k_\chi^-$, then some class contains at
least two vertices: otherwise $h=n$, contradicting $h<k_\chi^-<n$. Splitting
such a class into two nonempty subsets preserves independence. Repeating this
operation produces a proper coloring with exactly $k_\chi^-$ nonempty classes.
The size cap is preserved under splitting. Therefore
\[
  \{\chi(G_n)\le k_\chi^-\}\cap\mathcal A_n
  \subseteq
  \{\text{an exactly $k_\chi^-$, $(\alpha+1)$-bounded coloring exists}\}.
\]
Markov's inequality and (4.4a) now give
\begin{equation*}
\begin{aligned}
  \Prob{\chi(G_n)\le k_\chi^-}
  &\le
  \Prob{\mathcal A_n^c}
  +E_{n,k_\chi^-,\alpha+1}\\
  &\le
  \varepsilon_n^{\mathrm{cap}}
  +\varepsilon_n^{\mathrm{prof}}
  \longrightarrow0.
\end{aligned}
  \tag{4.5}
\end{equation*}
Equivalently,
\begin{equation*}
  \Prob{\chi(G_n)>k_\chi^-}\longrightarrow1.
  \tag{4.6}
\end{equation*}
The direction is strict, as required by the final event intersection.

Finally, (4.3b) also records the location precision:
\begin{equation*}
  |k_\chi^--r_+(n)|
  <\log n+2
  =o\!\left(\frac{n}{(\log n)^3}\right).
  \tag{4.7}
\end{equation*}
Thus both the probability statement and the location error hold with one set
of phase-independent thresholds along the full sequence of integers.
